    \begin{center}
        {\large \textbf{Вариант №22 \\
        }}
    \end{center}

    \textbf{\underline{Цель работы:} изучить основные возможности языка SQL для создания структуры базы данных. Научиться создавать базы данных, таблицы, связи, ограничения, а также создавать, изменять и удалять данные.
    } 

    \begin{enumerate}
        \item \hyperref[sec:1]{Составить SQL-запросы для создания структуры базы данных, полученной в результате лабораторной работы №1. Указать используемые типы данных, ограничения значений полей; для связей: действия с записями подчинённой таблицы при удалении и изменении соответствующей записи главной таблицы. }
        \item \hyperref[sec:2]{С помощью SQL-запросов выполнить добавление 3–4 записей в каждую таблицу, изменение и удаление нескольких записей.}
    \end{enumerate}